\section{源码等价数学实现}

本章只转写当前源码中的赋值、分支、截断与迭代，不使用教材化通式。符号名称尽量保持与
结构字段和局部变量一致；每组公式后的路径与行号是本章的核验边界。

\subsection{Holland 风速实现}

\srcpath{holland\_wind\_ms} 直接调用
\srcpath{holland\_wind\_ms\_impl}。对轨迹点 \texttt{storm} 和站点
$(\texttt{site\_lat},\texttt{site\_lon})$，代码依次计算
\begin{align}
 \texttt{radius\_km}&=\max\!\left(1,
   \operatorname{haversine\_km}(\texttt{storm.latitude},\texttt{storm.longitude},
   \texttt{site\_lat},\texttt{site\_lon})\right),\\
 \texttt{rmw}&=\max(1,\texttt{storm.rmw\_km}),\\
 \texttt{b}&=\max(0.5,\texttt{storm.holland\_b}),\\
 \texttt{ratio}&=(\texttt{rmw}/\texttt{radius\_km})^{\texttt b},\\
 \texttt{pressure}&=
   \frac{\texttt b\,\texttt{storm.delta\_p\_hpa}\,100\,
   \texttt{ratio}\,e^{-\texttt{ratio}}}{\texttt{kAirDensity}},\\
 \texttt{r\_m}&=1000\,\texttt{radius\_km},\qquad
 \texttt{tempc}=\texttt{storm.fc}\,\texttt{r\_m}/2,\\
 \texttt{gradient}&=\sqrt{\max(0,\texttt{pressure}+\texttt{tempc}^{2})}-\texttt{tempc},\\
 \texttt{v10}&=\texttt{kGustFactor}\,\texttt{gradient}
 \left(\frac{\texttt{kAnemometerHeightM}}{\texttt{kGradientHeightM}}\right)^{\texttt{kPowerLawExponent}}.
\end{align}
返回值采用严格分支：当 \texttt{v10}<0.01 时返回 0，否则返回
$\max(0,\texttt{v10})$。当前实现没有计算或返回气压廓线 $p(r)$，也没有把平移速度矢量叠加到
\texttt{gradient}，因此本手册不写这两项。依据为
\srcpath{src/scenario\_generation/typhoon\_resilience.cpp:holland\_wind\_ms\_impl} 和
\srcpath{src/scenario\_generation/typhoon\_resilience.cpp:holland\_wind\_ms}。

\subsection{轨迹点 RMW 与 Holland B 生成}

\srcpath{build\_track} 对每个轨迹点按以下分支赋值
\begin{align}
 \texttt{p.rmw\_km}&=\begin{cases}
 \texttt{opts.initial\_rmw\_km},&\texttt{opts.initial\_rmw\_km}>0,\\
 \operatorname{rmw\_from\_delta\_p}(\texttt{delta\_p}),&\text{否则},
 \end{cases}\\
 \texttt{p.holland\_b}&=\operatorname{holland\_b}(\texttt{lat},\texttt{p.rmw\_km}),
\end{align}
其中两个生成函数原样为
\begin{align}
 \operatorname{rmw\_from\_delta\_p}(\Delta p)
 &=\exp\!\left(3.858-7.7\times10^{-5}\,\Delta p^{2}\right),\\
 \operatorname{holland\_b}(\texttt{lat},\texttt{rmw\_km})
 &=\max\!\left(0.8,\;1.881093-0.010917\,\texttt{lat}
   -0.005567\,\texttt{rmw\_km}\right).
\end{align}
依据为 \srcpath{src/scenario\_generation/typhoon\_resilience.cpp:rmw\_from\_delta\_p} 和
\srcpath{src/scenario\_generation/typhoon\_resilience.cpp:holland\_b}。

\subsection{降雨实现}

令 \texttt{radius} 为台风中心到站点的 Haversine 距离与
\texttt{kRainRadiusFloorKm} 的较大值，并执行
\begin{align}
 q&=\operatorname{clamp}(\texttt{storm.rmw\_km}/\texttt{radius},0,
       \texttt{kRainRatioMax}),\\
 rr&=\max(0,-5.5+110q-390q^2+550q^3-250q^4),\\
 \texttt{dpdt}&=\begin{cases}
 \texttt{previous.delta\_p\_hpa}-\texttt{storm.delta\_p\_hpa},&\texttt{previous}\ne\varnothing,\\
 1,&\texttt{previous}=\varnothing,
 \end{cases}\\
 k&=\max(1,0.0319\,\texttt{storm.delta\_p\_hpa}-0.0395),\\
 k_1&=\max(1,1-\texttt{dpdt}/100),\\
 s&=\left\lfloor\operatorname{bearing\_sector\_deg}/45\right\rfloor\bmod 8,\\
 R&=\max(0,k k_1 rr\,m_s).
\end{align}
$m_s$ 完全按 \srcpath{rain\_sector\_multiplier} 的八分区和
\texttt{translation\_speed\_kmph}/3.6 的 $<4$、$>8$ 分支选取；不是连续拟合函数。
依据为 \srcpath{src/scenario\_generation/typhoon\_resilience.cpp:bearing\_sector\_deg}。

\subsection{架空段故障概率}

源码先计算雨致等效风速：
\begin{align}
 v_{10\min}&=\texttt{wind\_ms}/1.42,\\
 f_1&=e^{0.006462v_{10\min}}-1.2486e^{-0.2769v_{10\min}},\\
 f_2&=\begin{cases}0.09376\,\texttt{rain\_mm\_hr}^{0.7087},&\texttt{rain\_mm\_hr}>0,\\0,&\text{否则},\end{cases}\\
 v_{eq}&=\max(0,(v_{10\min}+\max(0,f_1f_2))1.42).
\end{align}
仅当 $v_{eq}>\texttt{design\_wind\_ms}$ 时，才令
\begin{align}
 z&=\frac{\ln v_{eq}-\ln\texttt{lognormal\_median\_wind\_ms}}
          {\texttt{lognormal\_sigma}},\\
 \texttt{severity}&=\operatorname{clamp}\!\left(\tfrac12
       \operatorname{erfc}(-z/\sqrt2),0,1\right),\\
 \ell&=\operatorname{clamp}(\texttt{seg.length\_km},0,5),\\
 p_{wind}&=\operatorname{clamp}\!\left(1-e^{-0.015\,
       \texttt{severity}\,\ell\,\max(0,\texttt{dt\_hr})},0,1\right);
\end{align}
否则 $p_{wind}=0$。若对数正态函数任一输入不为正，\texttt{severity} 的底层函数返回 0。

另一个段暴露项不是 logistic。源码计算
\begin{align}
 r_v&=\begin{cases}\texttt{wind\_ms}/\texttt{design\_wind\_ms},&
 \texttt{design\_wind\_ms}>0,\\0,&\text{否则},\end{cases}\\
 r_R&=\begin{cases}\texttt{rain\_mm\_hr}/\texttt{design\_rain\_mm\_hr},&
 \texttt{design\_rain\_mm\_hr}>0,\\0,&\text{否则},\end{cases}\\
 \lambda&=\texttt{seg.length\_km}^{2}
 e^{\texttt{seg\_a\_wind}r_v+\texttt{seg\_b\_rain}r_R+\texttt{seg\_c\_bias}},\\
 p_{seg}&=\begin{cases}
 \operatorname{clamp}(1-e^{-\lambda\max(0,\texttt{dt\_hr})},0,1),&
 \texttt{wind\_ms}>\texttt{design\_wind\_ms}\ \lor\
 \texttt{rain\_mm\_hr}>\texttt{design\_rain\_mm\_hr},\\
 0,&\text{否则},
 \end{cases}\\
 p_{over}&=\operatorname{clamp}(1-(1-p_{wind})(1-p_{seg}),0,1).
\end{align}
依据为 \srcpath{src/scenario\_generation/typhoon\_resilience.cpp:normal\_cdf}。

\subsection{电缆段故障概率与资产分支}

电缆分支逐项执行
\begin{align}
 CN&=\operatorname{clamp}(\texttt{curve\_number\_cn},1,99),&
 S&=25400/CN-254,& I_a&=0.2S,\\
 \texttt{runoff}&=\begin{cases}
 (R-I_a)^2/\max(10^{-9},R-I_a+S),&R>I_a,\\0,&R\le I_a,
 \end{cases}\\
 \texttt{infiltration}&=\max(0,R-\texttt{runoff}),&
 q&=\texttt{runoff}/1000/3600,\\
 h_t&=\left(\frac{q\,\texttt{manning\_n}}
 {\sqrt{\max(10^{-9},\texttt{slope\_s0})}}\right)^{0.6},\\
 \theta&=\min\!\left(\texttt{porosity\_phi},
 \texttt{initial\_moisture\_theta}+
 \frac{\texttt{infiltration}}
 {\max(10^{-9},1000\,\texttt{soil\_depth\_z\_m})}\right),\\
 p_{cable}&=\operatorname{clamp}\!\left(
 \frac{\texttt{base\_cable\_failure\_p}}{10}
 e^{130h_t+8\theta/\max(10^{-9},\texttt{porosity\_phi})},0,1\right).
\end{align}
\texttt{Cable} 直接返回 $p_{cable}$；\texttt{Mixed} 返回
$\operatorname{clamp}(1-(1-p_{over})(1-p_{cable}),0,1)$；其他类型返回 $p_{over}$。
依据为 \srcpath{src/scenario\_generation/typhoon\_resilience.cpp:cable\_segment\_probability}。

\subsection{支路聚合、抽样和修复}

每一时步内，支路所有段的生存概率连乘：
\begin{equation}
 p_{branch,t}=\operatorname{clamp}\!\left(1-\prod_{s\in branch}
 (1-\operatorname{clamp}(p_{s,t},0,1)),0,1\right).
\end{equation}
每条尚未故障的支路在每个时步独立抽取 $U\sim U[0,1]$；首次满足
$U<p_{branch,t}$ 即记录故障时刻。\fld{min\_fault\_probability} 不参与这段故障过滤。
基础修复时间优先取 AC 的 \fld{mttr\_hr} 或 DC 的 \fld{mttr\_hours}（仅正值），否则取
\fld{default\_repair\_time\_hr}，随后执行
\begin{align}
 T_r&\leftarrow T_r\left(1+\texttt{repair\_time\_wind\_factor}
 \frac{\max(0,\texttt{peak\_wind\_ms}-25)}{25}\right),\\
 T_r&\leftarrow\operatorname{clamp}(T_r,
 \max(10^{-6},\texttt{time\_step\_hr}),\texttt{max\_repair\_time\_hr}).
\end{align}
依据为 \srcpath{src/scenario\_generation/typhoon\_resilience.cpp:generate\_typhoon\_fault\_sequence}。

启用分阶段修复时，准备时刻为“最后故障时刻加非负延迟”，每槽时间为
$\max(\texttt{time\_step\_hr},\texttt{repair\_crew\_time\_per\_branch\_hr})$。
精确分支按恢复后的供电增量加以下原样保留的破同分项选择最大者：
\begin{equation}
 \texttt{score}=\texttt{impact}+10^{-3}\texttt{peak\_probability}
 +10^{-5}\texttt{peak\_wind\_ms}+10^{-6}\texttt{branch\_length\_km}.
\end{equation}
当故障数超过 \texttt{kMaxExactStagedRepairFaults} 时，改用源距离、故障时间、峰值概率、
支路类型、稳定 ID 的顺序排序，并置
\fld{used\_approximate\_repair\_order}=true。依据为
\srcpath{src/scenario\_generation/typhoon\_resilience.cpp:apply\_staged\_post\_disaster\_repairs}。

\subsection{交通影响递推}

每条交通边、每个采样步执行
\begin{align}
 W_{t+1}&=\operatorname{clamp}\!\left(W_t+
 (\texttt{rainfall\_runoff\_coefficient}\,
 \texttt{runoff\_multiplier}\,R_t-\texttt{drainage})\Delta t,
 0,\texttt{maximum\_surface\_water\_mm}\right),\\
 f_v&=\operatorname{clamp}\!\left(
 \frac{aW_{t+1}^{2}+bW_{t+1}+c}{\max(\texttt{kTol},c)},
 \texttt{minimum\_open\_speed\_factor},1\right),\\
 f_w&=1-\operatorname{clamp}\!\left(
 \frac{V-\texttt{wind\_capacity\_reduction\_start\_ms}}
 {\max(\texttt{kTol},\texttt{wind\_capacity\_reduction\_full\_ms}-
 \texttt{wind\_capacity\_reduction\_start\_ms})},0,1\right)
 (1-\operatorname{clamp}(\texttt{minimum\_wind\_capacity\_factor},0,1)),\\
 f_c&=\operatorname{clamp}\!\left(
 f_v^{\max(0,\texttt{capacity\_speed\_exponent})}f_w,
 \texttt{minimum\_open\_capacity\_factor},1\right).
\end{align}
当 $V\le\texttt{wind\_capacity\_reduction\_start\_ms}$ 时源码直接令 $f_w=1$。
边可用当且仅当原边可用、$W_{t+1}<\texttt{flood\_closure\_depth\_mm}$，且未启用风闭锁或
$V<\texttt{wind\_closure\_ms}$。时耗更新为
$\texttt{base\_time\_hr}/\max(\texttt{minimum\_open\_speed\_factor},f_v)$；容量在可用时为
$\texttt{base\_capacity}f_c$，否则为 0。依据为
\srcpath{src/scenario\_generation/typhoon\_traffic\_impact.cpp:wind\_capacity\_factor} 和
\srcpath{src/scenario\_generation/typhoon\_traffic\_impact.cpp:apply\_typhoon\_traffic\_impact}。

\subsection{候选距离与 medoid 更新}

对已经由 \srcpath{scaled\_features} 处理的连续向量 $a,b$，源码距离为
\begin{equation}
 d=\texttt{continuous\_weight}\sqrt{\sum_i(a_i-b_i)^2}
 +\texttt{outage\_hamming\_weight}\sum_i[\sigma_i^a\ne\sigma_i^b]
 +\texttt{regime\_weight}[r_a\ne r_b].
\end{equation}
最后一项仅在混合方法、两个标签均非空且标签不同的情况下取 1；长度不足的停运签名以 0
补齐。分配步把候选分给严格距离最小的 medoid；非冻结簇的更新步在该簇候选中最小化
\begin{equation}
 \sum_{j:\,assignment_j=c}\texttt{candidates}[j].\texttt{probability}\ d(i,j).
\end{equation}
迭代上限是 $\max(1,\texttt{max\_iterations})$，无变化即提前停止。初始化的抽样权重是
$\max(10^{-12},\pi_i)d_{nearest,i}^{2}$，并非普通无权 k-medoids。聚类概率先累加成员概率，
当总和大于 $10^{-12}$ 时再除以总和。依据为
\srcpath{src/scenario\_generation/scenario\_generation.cpp:candidate\_distance}。

\subsection{风险分数和尾部锚点的实际语义}

各风险源列先以中位数居中，以
$\max(10^{-9},Q_{75}-Q_{25})$ 缩放。对候选的缩放值集合 $v_k$，源码赋值
\begin{equation}
 \texttt{risk\_score}=\max_k v_k
 +0.25\sqrt{\sum_{v_k>1}(v_k-1)^2}
 +0.15\sqrt{\frac{1}{K}\sum_k(v_k-\bar v)^2}.
\end{equation}
全局尾部和单源尾部分别按 \fld{source\_tail\_quantile} 分位数判定；至少两个源达到
\fld{coupling\_quantile} 时增加耦合尾标签。锚点集合还包含有故障的弹性候选以及每个
\fld{regime\_label} 中风险最高的候选；若超过 $k$，优先保留有故障候选，再按
\fld{risk\_score} 降序截断到 $k$。因此 \fld{tail\_fraction} 并未参与当前锚点选择，不能描述为
“固定保留最高 5\%”。依据为 \srcpath{src/scenario\_generation/scenario\_generation.cpp:raw\_risk\_source\_scores}。

\subsection{覆盖边界}

本章覆盖当前实现中可独立核验的台风风雨、段/支路故障、修复、交通递推和聚类公式。
常规、可靠性、弹性候选的概率在各生成分支中直接按有效候选数设为等概率；可靠性与弹性分支
分别见 \srcpath{src/scenario\_generation/scenario\_generation.cpp:base\_load\_components（lambda）}、
\srcpath{src/scenario\_generation/scenario\_generation.cpp:generate\_resilience\_scenarios}。本章没有把配置字段名称转化为代码不存在的
随机分布、物理约束或跨族概率模型。
